\chapter{Introduzione}
\label{ch:introduzione}

% generazione di test per progetti python con large language models
% quanto sono in grado di coprire il codice effettivamente (coverage)
% possibile misura di quanto sono rieseguite le linee (contare quante volte viene eseguita la stessa linea)
% quante linee sono replicate (limitato ai file di test, per esempio con questo tool https://github.com/platisd/duplicate-code-detection-tool)


% --- da qui: bozza da rielaborare ---

Il testing di unità è una componente critica del processo di sviluppo del
software: una suite di test di buona qualità permette di individuare le
incongruenze tra le specifiche di un sistema e la sua implementazione, e di
catturare le regressioni introdotte durante l'evoluzione del codice. Scrivere
buoni test è tuttavia un'attività costosa e ripetitiva, che nella pratica viene
spesso trascurata in tutto o in parte.

I modelli linguistici di grandi dimensioni (LLM) hanno mostrato capacità
notevoli nella generazione di codice, e la generazione automatica di unit test
è una delle applicazioni più promettenti. Valutare la qualità dei test così
prodotti è però tutt'altro che immediato: un test può essere sintatticamente
valido ma non eseguibile, può eseguire il codice senza verificarne davvero il
comportamento, oppure può raggiungere una copertura elevata pur non essendo in
grado di individuare alcun difetto.

\section{Obiettivo della tesi}
% TODO: rifinire con la relatrice.

Questa tesi studia \emph{con quali metriche} si può valutare la qualità degli
unit test generati da un LLM, e cosa ciascuna metrica riesce a cogliere. A
questo scopo vengono generati test per un insieme di funzioni Python reali
utilizzando tre modelli della stessa famiglia ma di dimensione crescente, a
parità di prompt, di dataset e di parametri di generazione; i test prodotti
vengono poi eseguiti e misurati secondo più dimensioni: esito dell'esecuzione,
copertura del codice, ridondanza fra i test e aderenza al formato richiesto.

Il confronto fra modelli di taglia diversa a condizioni identiche permette di
osservare quali aspetti della qualità migliorano con la dimensione del modello
e quali no; l'uso congiunto di più metriche permette di mettere in luce
fenomeni che una singola misura nasconderebbe.

\section{Struttura del documento}
% TODO: completare quando i capitoli saranno stabili.

Il Capitolo~\ref{ch:stato_arte} inquadra lo stato dell'arte sulla generazione
automatica di test e sulle metriche usate per valutarla. Il
Capitolo~\ref{ch:core} descrive il sistema realizzato. Il
Capitolo~\ref{ch:risultati} presenta gli esperimenti e i risultati ottenuti.
Seguono le conclusioni e gli sviluppi futuri.
